% ======================================================================
% capitulos/capitulo4.tex
% ----------------------------------------------------------------------
% ENTREGABLE 4 — MODELADO DE PROCESOS AS-IS / TO-BE
% Empresa caso de estudio: WISP Ingeniería en Redes y Telecomunicaciones
% Proceso modelado: Ciclo de vida del cliente (captación -> activación)
% Contenido trasladado desde Entregrable4_Procesos_ASIS_TOBE_WISP.docx
% ======================================================================
\chapter{Modelado de Procesos AS-IS / TO-BE}
\label{chap:procesos}

% ----------------------------------------------------------------------
% Descripción del producto (introducción del capítulo)
% ----------------------------------------------------------------------
El presente documento contiene la representación estructurada de los procesos
organizacionales actuales (AS-IS) de la empresa WISP y su rediseño propuesto
(TO-BE). Permite visualizar las ineficiencias, redundancias y puntos críticos del
modo de operar vigente, y proponer mejoras mediante la automatización y la
reestructuración apoyada en TIC.

El proceso seleccionado para el modelado es el \textbf{ciclo de vida del cliente},
que abarca desde que un vendedor capta a un prospecto en campo hasta que el cliente
queda activo con servicio de internet. Se eligió este proceso \textit{end-to-end}
porque integra en un solo flujo los tres problemas operativos críticos de WISP
(captación, instalación y activación en red), lo que permite cuantificar de forma
agregada la mejora en tiempo, costo y calidad del servicio.

\begin{upeubox}{Fundamento metodológico: notación BPMN 2.0}
El modelado utiliza la notación BPMN 2.0 (\textit{Business Process Model and
Notation}, estándar OMG). Los diagramas emplean piscinas y carriles
(\textit{pools/lanes}) para representar los responsables, tareas (rectángulos
redondeados), compuertas de decisión (rombos), eventos de inicio y fin (círculos) y
flujos de secuencia (flechas), garantizando una representación normalizada e
interpretable del proceso.
\end{upeubox}

% ======================================================================
\section{Proceso Actual (AS-IS)}
\label{sec:asis}

\subsection{Descripción narrativa del proceso}

Actualmente, el ciclo de vida del cliente en WISP se desarrolla de la siguiente
manera: un vendedor de campo capta a un prospecto y le explica el servicio, en una
conversación que dura entre 5 y 10 minutos. Los datos del prospecto se anotan en
papel o se registran informalmente por WhatsApp. Luego, el vendedor debe avisar
manualmente al administrador (\emph{«tengo un cliente para instalar»}), ya que no
existe ningún mecanismo de notificación automática; este aviso y la transferencia de
datos pueden demorar desde 1 minuto hasta 2 horas, según la disponibilidad del
administrador.

El administrador recibe los datos, evalúa si la instalación procede según la
cobertura de la zona, y sube manualmente la orden al sistema de órdenes. El técnico
visualiza la orden y realiza la instalación, generalmente de un día para otro (si el
cliente solicita hoy, se instala al día siguiente); como excepción, si el técnico ya
se encuentra en la zona y capta directamente al cliente, puede instalar en el
momento. Tras la instalación, el técnico avisa nuevamente, y el administrador debe
dar de alta manualmente al cliente como usuario PPPoE en el MikroTik de la zona
correspondiente, tarea que toma alrededor de 15 minutos. Recién entonces el cliente
queda activo con servicio.

\subsection{Diagrama BPMN — Proceso Actual (AS-IS)}

\begin{center}
\includegraphics[width=\textwidth]{bpmn_asis.png}
\captionof{figure}{Diagrama BPMN del ciclo de vida del cliente actual (AS-IS). En
rojo, las tareas manuales críticas que generan demoras y pérdidas.}
\label{fig:bpmn-asis}
\end{center}

\subsection{Indicadores actuales (AS-IS)}

A partir de los tiempos reales observados en la operación de WISP, se establecen los
siguientes indicadores del proceso actual:

\begin{center}
\captionof{table}{Indicadores del proceso actual (AS-IS).}
\label{tab:ind-asis}
\begin{tabularx}{\textwidth}{L{4.3cm} L{2.6cm} X}
\toprule
\textbf{Indicador} & \textbf{Valor actual} & \textbf{Observación} \\
\midrule
Tiempo de captación & 5--10 min & Conversación con el prospecto (humano). \\
Tiempo de transferencia al admin & 1 min -- 2 h &
Aviso manual, sin notificación automática. Cuello de botella. \\
Alta manual en MikroTik & $\sim$15 min & El admin entra manualmente a cada MikroTik. \\
Lead time total (captación \textrightarrow{} activación) & $\sim$27 h ($\sim$1.1 días) &
Incluye espera de instalación de un día para otro. \\
Tiempo de trabajo efectivo & $\sim$198 min & Suma de tareas activas (sin esperas). \\
Clientes perdidos en el proceso & 1--2 por ciclo & Prospectos captados que nunca se activan. \\
Trazabilidad de la instalación & 0\% & Evidencias por WhatsApp, sin vínculo al cliente. \\
\bottomrule
\end{tabularx}
\end{center}

\subsection{Problemas detectados}

\begin{itemize}[leftmargin=1.4em]
  \item \textbf{Dependencia del aviso manual:} el proceso se detiene hasta que el
  vendedor logra contactar al administrador; no hay notificaciones automáticas.
  \item \textbf{Pérdida y dispersión de datos:} la información se anota en papel o
  WhatsApp, lo que genera errores de transcripción y prospectos perdidos.
  \item \textbf{Doble digitación:} los datos se registran informalmente y luego el
  admin los reingresa al sistema de órdenes.
  \item \textbf{Alta manual en MikroTik:} el administrador debe ingresar a cada
  router para activar al cliente, tarea repetitiva y propensa a error en 13 zonas.
  \item \textbf{Sin trazabilidad técnica:} las evidencias de instalación (caja NAP,
  borne, equipos) se envían por WhatsApp y no quedan vinculadas al cliente.
  \item \textbf{Cuello de botella en el administrador:} toda la operación depende de
  una sola persona disponible.
\end{itemize}

% ======================================================================
\section{Proceso Propuesto (TO-BE)}
\label{sec:tobe}

\subsection{Rediseño del flujo}

El proceso rediseñado incorpora el sistema de gestión integral como un actor activo
(carril «Sistema») que automatiza las transferencias de información y elimina los
avisos manuales. El vendedor registra al prospecto directamente en la aplicación
móvil, en el momento y lugar de la captación. El sistema valida automáticamente la
cobertura de la zona y notifica de inmediato al administrador, sin que el vendedor
tenga que contactarlo.

El administrador valida el prospecto con un solo clic (los datos ya están
estructurados y disponibles), y el sistema genera automáticamente la orden de
instalación, notificando al técnico en su aplicación. El técnico realiza la
instalación y registra la evidencia (fotos de caja NAP, borne, equipos, fachada)
directamente en la app, vinculada al cliente. Al cerrar la orden, el sistema da de
alta automáticamente al cliente como usuario PPPoE en el MikroTik correspondiente
mediante la API de RouterOS, dejándolo activo de forma inmediata.

\subsection{Diagrama BPMN — Proceso Propuesto (TO-BE)}

\begin{center}
\includegraphics[width=\textwidth]{bpmn_tobe.png}
\captionof{figure}{Diagrama BPMN del ciclo de vida del cliente propuesto (TO-BE). El
carril «Sistema» automatiza validaciones, notificaciones y el alta en MikroTik.}
\label{fig:bpmn-tobe}
\end{center}

\subsection{Automatizaciones propuestas}

El rediseño introduce las siguientes automatizaciones apoyadas en TIC:

\begin{center}
\captionof{table}{Automatizaciones introducidas por el rediseño.}
\label{tab:automatizaciones}
\begin{tabularx}{\textwidth}{L{3.8cm} X L{3.8cm}}
\toprule
\textbf{Automatización} & \textbf{Reemplaza a} & \textbf{Tecnología} \\
\midrule
Registro digital en campo & Anotación en papel/WhatsApp & App móvil (rol vendedor) \\
Notificación automática al admin & Aviso manual (1 min--2 h) & Notificaciones push / backend \\
Validación de cobertura & Verificación mental del admin & Sistema (zonas georreferenciadas) \\
Generación de orden & Subida manual a órdenes & Backend Django \\
Registro de evidencia técnica & Fotos sueltas por WhatsApp & App móvil (rol técnico) \\
Alta automática en MikroTik & Alta manual ($\sim$15 min) & API RouterOS $+$ WireGuard \\
\bottomrule
\end{tabularx}
\end{center}

\subsection{Nuevos indicadores (TO-BE)}

\begin{center}
\captionof{table}{Indicadores del proceso propuesto (TO-BE).}
\label{tab:ind-tobe}
\begin{tabularx}{\textwidth}{L{4.3cm} L{3.2cm} X}
\toprule
\textbf{Indicador} & \textbf{Valor propuesto} & \textbf{Mejora} \\
\midrule
Tiempo de transferencia al admin & Instantáneo ($\sim$0.1 min) & Notificación automática \\
Alta en MikroTik & $\sim$0.5 min (automática) & Vía API RouterOS \\
Lead time total & $\sim$14 h ($\sim$0.6 días) & Mejor coordinación \\
Tiempo de trabajo efectivo & $\sim$111 min & Menos pasos manuales \\
Tiempo de admin por cliente & $\sim$5.5 min & Antes $\sim$86 min \\
Trazabilidad de instalación & 100\% & Evidencia vinculada \\
Clientes perdidos & $\approx$ 0 & Registro inmediato \\
\bottomrule
\end{tabularx}
\end{center}

\begin{upeubox}{Mejora estimada en eficiencia}
El rediseño reduce el tiempo total del ciclo (\textit{lead time}) en aproximadamente
un \textbf{49\%}, el tiempo de trabajo efectivo en un \textbf{44\%}, y el tiempo de
gestión del administrador por cliente en un \textbf{94\%}. El impacto más
significativo se da en la transferencia de información (de hasta 2 horas a
instantánea, una reducción del 96\%) y en el alta en MikroTik (de 15 minutos manuales
a medio minuto automático, un 97\% menos).
\end{upeubox}

% ======================================================================
\section{Análisis Comparativo}
\label{sec:comparativo}

Se comparan cuantitativamente los procesos AS-IS y TO-BE en las tres dimensiones
exigidas: tiempo, costo y calidad.

\subsection{Reducción de tiempos}

\begin{center}
\captionof{table}{Comparación de tiempos entre AS-IS y TO-BE.}
\label{tab:comp-tiempos}
\begin{tabularx}{\textwidth}{L{3.8cm} X X C{1.9cm}}
\toprule
\textbf{Métrica de tiempo} & \textbf{AS-IS} & \textbf{TO-BE} & \textbf{Reducción} \\
\midrule
Transferencia al admin & 1 min--2 h & Instantáneo & 96\% \\
Alta en MikroTik & $\sim$15 min & $\sim$0.5 min & 97\% \\
Trabajo efectivo total & $\sim$198 min & $\sim$111 min & 44\% \\
Lead time (ciclo completo) & $\sim$27 h & $\sim$14 h & 49\% \\
Gestión admin por cliente & $\sim$86 min & $\sim$5.5 min & 94\% \\
\bottomrule
\end{tabularx}
\end{center}

\subsection{Reducción de costos}

La reducción de tiempos se traduce en ahorro de costos operativos, principalmente por
la liberación de tiempo del administrador y la reducción de clientes perdidos:

\begin{center}
\captionof{table}{Comparación de costos entre AS-IS y TO-BE.}
\label{tab:comp-costos}
\begin{tabularx}{\textwidth}{L{4cm} X X X}
\toprule
\textbf{Concepto} & \textbf{AS-IS} & \textbf{TO-BE} & \textbf{Impacto} \\
\midrule
Tiempo admin/cliente valorizado & $\sim$86 min & $\sim$5.5 min & $-$94\% \\
Clientes perdidos por ciclo & 1--2 & $\approx$ 0 & Ingreso recuperado \\
Retrabajo (doble digitación) & Sí & Eliminado & $-$100\% \\
\bottomrule
\end{tabularx}
\end{center}

\textit{Nota:} la valorización económica detallada de estos ahorros se desarrolla en
el Business Case (beneficio de ahorro administrativo y recuperación por menor pérdida
de clientes).

\subsection{Mejora en calidad del servicio}

\begin{center}
\captionof{table}{Comparación de la calidad del servicio entre AS-IS y TO-BE.}
\label{tab:comp-calidad}
\begin{tabularx}{\textwidth}{L{4cm} X X}
\toprule
\textbf{Dimensión de calidad} & \textbf{AS-IS} & \textbf{TO-BE} \\
\midrule
Trazabilidad del cliente & 0\% (WhatsApp/papel) & 100\% en el sistema \\
Historial técnico (NAP, borne, ONU) & Inexistente & Registrado y vinculado \\
Errores de transcripción & Frecuentes & Minimizados (registro directo) \\
Consistencia de datos & Datos dispersos & Fuente única de verdad \\
Experiencia del cliente & Activación variable & Activación más rápida \\
Dependencia de una persona & Total (admin) & Distribuida (sistema) \\
\bottomrule
\end{tabularx}
\end{center}

\begin{upeubox}{Conclusión del análisis comparativo}
El rediseño del proceso (TO-BE) mejora simultáneamente las tres dimensiones: reduce
el tiempo del ciclo en $\sim$49\% y la carga del administrador en $\sim$94\%, elimina
el retrabajo y la pérdida de clientes, y eleva la calidad al garantizar trazabilidad
total, historial técnico y una fuente única de datos. La mejora no es incremental sino
estructural: se pasa de un proceso dependiente de avisos manuales y personas, a un
proceso automatizado, trazable y escalable. Este rediseño constituye la justificación
operativa directa del sistema propuesto.
\end{upeubox}
